\Needspace{14\baselineskip}
\subsection{Total variation of the atomic data}
\label{sec:curve-atomic-topology}

The trace calculation extends to arbitrary ideal envelopes on $Y$. Their atomic curvature determines them, and taking a maximum replaces two atomic measures by their common part. The ordered mass formula can therefore compare envelopes whose poles are unrelated. For this subsection, fix a smooth positive form $\beta$ on $Y$, put $A=\int_Y\beta$, and choose a nonpositive Green kernel for $\eta=\beta/A$. All atomic measures below have countable support and finite total mass.

For a positive atomic measure $D$ with $D(Y)\leq A$, set
\[
 w_D(x)=\int_YG(x,p)\,\mathrm{d}D(p),\qquad
 q_D=P_\beta[w_D]_{\mathcal I}.
\]
The Green construction gives a genuine quasi-plurisubharmonic function
with $\ddc w_D=D-D(Y)\eta$.

\begin{lemma}\label{lem:curve-atomic-ideal-model}
The function $w_D$ is $\beta$-plurisubharmonic, and
\[
 q_D\sim w_D,\qquad
 \int_Y\mathrm{MA}_\beta(q_D)
 =\int_Y\mathrm{MA}_\beta(w_D)=A-D(Y).
\]
The map $D\mapsto q_D$ is a bijection onto all normalized ideal envelopes.
The positive-mass locus corresponds to $D(Y)<A$. The measure $D$ is the atomic part of
$\beta+\ddc q_D$.
\end{lemma}
\begin{proof}
One has
\[
 \beta+\ddc w_D=D+(A-D(Y))\eta\geq0.
\]
Every atom of $D$ is a pole of $w_D$. The smooth measure $\eta$ charges
no polar set, including any additional poles caused by accumulation of
the atoms. Hence $\mathrm{MA}_\beta(w_D)=(A-D(Y))\eta$.

Put $q=q_D$. Ideal projection preserves the point Lelong numbers. Applied to $\ddc w_D=D-D(Y)\eta$, \cref{lem:divisorial-potential-domination} gives $q\preceq w_D$. The envelope inequality $q\geq w_D-\sup_Yw_D$ proves $q\sim w_D$. By \cref{cor:polar-order-maximum}, the polar part of $\beta+\ddc q$ is therefore also $D$. Stokes gives its non-pluripolar mass $A-D(Y)$.

For any $u\in\PSH(Y,\beta)$, its point Lelong numbers give an atomic
measure $D_u$ of mass at most $A$. On a smooth curve these numbers
determine all multiplier ideals, by \cref{thm:local-i-inputs}.
Thus $u\sim_{\mathcal I}w_{D_u}$, so
$P_\beta[u]_{\mathcal I}=q_{D_u}$.
If $D_u(Y)=A$, its atoms exhaust the full curvature measure and its
non-pluripolar mass is zero. This proves surjectivity and the description of the positive-mass locus. Distinct atomic measures have
distinct Lelong data, proving injectivity.
\end{proof}

\Needspace{12\baselineskip}
\begin{theorem}\label{thm:curve-atomic-source-metric}\index{Trace!atomic metric}
For positive atomic measures $D,E$ with masses less than $A$,
\[
 \frac1{96}\|D-E\|_{\mathrm{TV}}
 \leq d_{S,\beta}(q_D,q_E)
 \leq\frac12\|D-E\|_{\mathrm{TV}}.
\]
Thus the singularity distance is a metric on these ideal envelopes and
induces exactly total-variation convergence of their atomic data.
\end{theorem}
\begin{proof}
Write $d_S=d_{S,\beta}$ and $a\lor b=\max(a,b)$. For positive-mass potentials $u,v\in\PSH(Y,\beta)$, besides the ordered-distance formula we use the inequality
\cite[Lemma~6.2.1]{SGPT}
\[
 d_S(u,v)\leq d_S(u,u\lor v)+d_S(v,u\lor v)
             \leq48d_S(u,v)
\]
on a curve. Here $48=3(1+1)2^{1+2}$ is the constant in that lemma for dimension one.

Put $z=\max(w_D,w_E)$. The lemma identifies the polar curvature measures of $w_D,w_E$ as $D,E$. By \eqref{eq:polar-maximum-measure}, the polar curvature of $z$ is $D\wedge E$, the atomic measure with the smaller coefficient at each point. Stokes therefore gives
\[
 \mathrm{MA}_\beta(z)(Y)=A-(D\wedge E)(Y).
\]
The ordered-distance formula yields
\[
 \begin{aligned}
 d_S(w_D,z)+d_S(w_E,z)
 &=\frac12\bigl(D(Y)+E(Y)-2(D\wedge E)(Y)\bigr)\\
 &=\frac12\|D-E\|_{\mathrm{TV}}.
 \end{aligned}
\]
Combining this equality with the two inequalities above proves both bounds. Ordinary equivalence with $q_D,q_E$ leaves the singularity distance unchanged. The lower bound and the lemma's injectivity make the pseudometric a metric on these representatives.
\end{proof}

\Needspace{12\baselineskip}
\begin{corollary}\label{cor:curve-ideal-source-completeness}
For $0<\varepsilon<A$, the normalized ideal envelopes satisfying
$\mathrm{MA}_\beta(q)(Y)\geq\varepsilon$ form a complete space for
$d_{S,\beta}$. Analytic ideal envelopes are dense in this space.
The space is not totally bounded.
\end{corollary}
\begin{proof}
By the lemma, the space corresponds to the positive atomic measures of
mass at most $A-\varepsilon$. This set is complete for total variation.
Indeed a Cauchy sequence has a positive finite-measure limit with the
same mass bound. The union of the sequence's countable supports is
countable, and total-variation convergence gives zero limit mass off
that union. The limit is therefore still atomic. The two metric
estimates transfer completeness to the ideal envelopes.

Finite truncation followed by rational approximation of each coefficient from below preserves the mass bound and converges in total variation. The resulting Green potentials have finitely many rational logarithmic poles; their ideal envelopes have the same ordinary types by the lemma. They are analytic, and the upper metric estimate proves density.

Finally fix $0<a<A-\varepsilon$. The measures $a\delta_p$, for distinct
$p\in Y$, have pairwise total-variation distance $2a$. Their ideal
envelopes have pairwise $d_S$ distance at least $a/48$. This infinite
separated family rules out total boundedness.
\end{proof}

This topology differs from the $L^1$ ideal metric of \cref{prop:ideal-lone-metric}: that metric is totally bounded. The sequence in \cref{prop:positive-mass-ideal-discontinuity} makes the distinction quantitative, since it converges in the ideal metric while remaining at a fixed positive $d_S$ distance from its limit.
